\documentclass[spanish, 12pt, a4paper]{article}

\usepackage{name}

\begin{document}

\maketitle
\thispagestyle{fancy}
\setcounter{page}{1}
\maketitle
\begin{multicols}{2}

\section{Modelo de Goodwin}

Consideramos el modelo de Goodwin para la resolución de la expresión genética dado por

\begin{align*}
    \frac{d m}{d t} & =\alpha_m g_R(p)-\beta_m m, \\
    \frac{d e}{d t} & =\alpha_e m-\beta_e e, \\
    \frac{d p}{d t} & =\alpha_p e-\beta_p p,
\end{align*}

donde $m$ es la concentración de mARN, que produce una enzima $e$, a una tasa $\alpha_m$ y se degrada a una tasa $\beta_m$. 
La enzima $e$ produce un producto $p$ a una tasa $\alpha_e$ y se degrada a una tasa $\beta_e$ y 
finalmente, el producto $p$ se degrada a una tasa $\beta_p$. Además, este modelo implementa una regulación controlada por la proteína,
con una función de la forma

\begin{equation}
    g_R(p)=\frac{a}{b+c p^h},
\end{equation}

donde $a$, $b$, $c$ son parámetros y $h$ es el exponente de Hill. Realizamos simulaciones variando el exponente de Hill 
manteniendo los demás parámetros constantes $a = b = c = 1$, $\alpha_m = \alpha_e = \alpha_p = 1$ y $\beta_m = \beta_e = \beta_p = 0.1$.

Los resultados para diferentes valores de $h$ se muestran en la Figura \ref{fig:goodwin_h}.

\begin{figure}[H]
    \centering
    \includegraphics[width=1\linewidth]{../goodwinOscilaciones.pdf}
    \caption{Simulación del modelo de Goodwin para distintos valores del exponente de Hill $h$.}
    \label{fig:goodwin_h}
\end{figure}

Se observa que para $h = 1$, las concentraciones tienden a un valor constante luego de pasar por un transitorio donde
se produce una subida y bajada de las concentraciones. 
Para $h = 2$, tienden a un valor constante mucho más bajo realizando oscilaciones amortiguadas, mientras que para
$h = 8$ y $h = 16$ se observan oscilaciones sostenidas en el tiempo.
Que se observen oscilaciones para un exponente de Hill tan alto como $h = 8$ es un resultado quizás forzado, ya que en la práctica, un exponente de Hill tan alto no es algo esperado en la regulación genética.

\begin{figure}[H]
    \centering
    \includegraphics[width=1\linewidth]{../goodwin0.9.pdf}
    \caption{Simulación del modelo de Goodwin para un valor de $\beta_m = \beta_e = \beta_p = 0.9$.}
    \label{fig:goodwin_degradacion}
\end{figure}

En una situación en la que se observan oscilaciones mediante este modelo, es importante analizar qué sucede al aumentar los valores de las degradaciones $\beta_m$, $\beta_e$ y $\beta_p$.
Por simplicidad los tomamos todos iguales. En la Figura \ref{fig:goodwin_degradacion} se observa que al haber aumentado las degradaciones con respecto al caso anterior ($\beta = 0.1$), las oscilaciones se amortiguan y tienden a desaparecer.


\section{Switch genético}

Estudiamos la dinámica de un sistema de dos genes con represión mutua dado por:

\begin{align*} 
    \frac{d m_1}{d t} & =\alpha_m g_R\left(p_2\right)-\beta_m m_1, \\ 
    \frac{d m_2}{d t} & =\alpha_m g_R\left(p_1\right)-\beta_m m_2, \\ 
    \frac{d p_1}{d t} & =\alpha_p m_1-\beta_p p_1, \\ 
    \frac{d p_2}{d t} & =\alpha_p m_2-\beta_p p_2,
\end{align*}

Usando la condición de que $\beta_m \gg \beta_p$ podemos reducir el sistema a solo dos variables.
Realizando la aproximación de esclavitúd suponemos que $m_1$ y $m_2$ se encuentran en estado estacionario, es decir, $\frac{d m_1}{d t} = \frac{d m_2}{d t} = 0$.
De esta forma, obtenemos que

\begin{align*}
    m_1 & = \frac{\alpha_m g_R(p_2)}{\beta_m}, \\
    m_2 & = \frac{\alpha_m g_R(p_1)}{\beta_m}.
\end{align*}

Reemplazando estas expresiones en las ecuaciones de $p_1$ y $p_2$ obtenemos

\begin{align*}
    \frac{d p_1}{d t} & = \alpha_p \frac{\alpha_m g_R(p_2)}{\beta_m} - \beta_p p_1, \\
    \frac{d p_2}{d t} & = \alpha_p \frac{\alpha_m g_R(p_1)}{\beta_m} - \beta_p p_2.
\end{align*}

Para estudiar las bifurcaciones que ocurren en el sistema dejamos fijos los parámetros
$\alpha_p = \alpha_m = \beta_m = 1$, $h = 2$, $a = c = 1$ y tomamos $\beta_p = 0.001$ para que se cumpla la aproximación.
Haciendo esto nos queda $b$ como parámetro de control. En las figuras \ref{fig:switch1}, \ref{fig:switch2} y \ref{fig:switch3} se observa
el cambio de comportamiento del sistema al variar $b$ en el diagrama de fases de las concentraciones $p_1$ y $p_2$.

\begin{figure}[H]
    \centering
    \includegraphics[width=1\linewidth]{../switch1.pdf}
    \caption{Diagrama de fases del sistema de dos genes con represión mutua simplificado, con $b = 40$.}
    \label{fig:switch1}
\end{figure}

\begin{figure}[H]
    \centering
    \includegraphics[width=1\linewidth]{../switch2.pdf}
    \caption{Diagrama de fases del sistema de dos genes con represión mutua simplificado, con $b = 60$.}
    \label{fig:switch2}
\end{figure}

\begin{figure}[H]
    \centering
    \includegraphics[width=1\linewidth]{../switch3.pdf}
    \caption{Diagrama de fases del sistema de dos genes con represión mutua simplificado, con $b = 80$.}
    \label{fig:switch3}
\end{figure}

En la Figura \ref{fig:switch1} se observa que para $b = 40$ el sistema tiene dos puntos fijos estables y uno inestable,
esto se corresponde con el hecho de que cuando uno de los genes está activo, suprime la
expresión del otro y viceversa.

En la Figura \ref{fig:switch2} se observa que para $b = 60$ las nulclinas de $p_1$ y $p_2$ se produce aproximadamente
la bifurcación, donde los puntos fijos estables desaparecen. Luego en la figura \ref{fig:switch3} se observa claramente el punto fijo
estable, este se corresponde con un estado donde ambos genes están activos.


\end{multicols}
\end{document}